% =====================================================================
\begin{thebibliography}{99}
\addcontentsline{toc}{chapter}{Bibliografía}
\markboth{Bibliografía}{}
\small
\def\UrlBreaks{\do\/\do-}

\bibitem{gqa}
J.~Ainslie, J.~Lee-Thorp, M.~de Jong, Y.~Zemlyanskiy, F.~Lebrón y S.~Sanghai.
\newblock GQA: Training generalized multi-query Transformer models from multi-head checkpoints.
\newblock En \emph{Proceedings of EMNLP}, 2023.
\newblock \url{https://arxiv.org/abs/2305.13245}

\bibitem{longformer}
I.~Beltagy, M.~E. Peters y A.~Cohan.
\newblock Longformer: The long-document Transformer.
\newblock Prepublicación, 2020.
\newblock \url{https://arxiv.org/abs/2004.05150}

\bibitem{pi}
S.~Chen, S.~Wong, L.~Chen e Y.~Tian.
\newblock Extending context window of large language models via position interpolation.
\newblock Prepublicación, 2023.
\newblock \url{https://arxiv.org/abs/2306.15595}

\bibitem{cmu11711}
Carnegie Mellon University.
\newblock 11-711: Advanced Natural Language Processing.
\newblock Otoño de 2024. Temario: mecanismos de atención, arquitecturas Transformer, preentrenamiento y multimodalidad.
\newblock \url{https://www.phontron.com/class/anlp-fall2024/}

\bibitem{cmu11868}
Carnegie Mellon University.
\newblock 11-868: Large Language Model Systems.
\newblock Primavera de 2025. Temario: entrenamiento distribuido, dispersión (MoE), compresión y sistemas de inferencia.
\newblock \url{https://llmsystem.github.io/llmsystem2025spring/}

\bibitem{deepseekv2}
DeepSeek-AI.
\newblock DeepSeek-V2: A strong, economical, and efficient mixture-of-experts language model.
\newblock Informe técnico, 2024.
\newblock \url{https://arxiv.org/abs/2405.04434}

\bibitem{deepseekv3}
DeepSeek-AI.
\newblock DeepSeek-V3 technical report.
\newblock Informe técnico, 2024.
\newblock \url{https://arxiv.org/abs/2412.19437}

\bibitem{encodec}
A.~Défossez, J.~Copet, G.~Synnaeve e Y.~Adi.
\newblock High fidelity neural audio compression.
\newblock Prepublicación, 2022.
\newblock \url{https://arxiv.org/abs/2210.13438}

\bibitem{vit}
A.~Dosovitskiy \emph{et al.}
\newblock An image is worth $16\times 16$ words: Transformers for image recognition at scale.
\newblock En \emph{International Conference on Learning Representations (ICLR)}, 2021.
\newblock \url{https://arxiv.org/abs/2010.11929}

\bibitem{switch}
W.~Fedus, B.~Zoph y N.~Shazeer.
\newblock Switch Transformers: Scaling to trillion parameter models with simple and efficient sparsity.
\newblock \emph{Journal of Machine Learning Research}, 23:1--40, 2022.
\newblock \url{https://arxiv.org/abs/2101.03961}

\bibitem{cs25}
S.~Feng, K.~Singh, M.~C. Frank y C.~D. Manning.
\newblock CS25: Transformers United, edición V6.
\newblock Stanford University, primavera. Sesiones de A.~Gu («On the Tradeoffs of State Space Models and Transformers»), N.~Tazi («The Ultra-Scale Talk: Scaling Training to Thousands of GPUs») y V.~Lin («From Language Models to Native Multimodal Intelligence»).
\newblock \url{https://web.stanford.edu/class/cs25/}

\bibitem{gemini15}
Gemini Team, Google.
\newblock Gemini 1.5: Unlocking multimodal understanding across millions of tokens of context.
\newblock Informe técnico, 2024.
\newblock \url{https://arxiv.org/abs/2403.05530}

\bibitem{llama3}
A.~Grattafiori \emph{et al.} (Llama Team, AI @ Meta).
\newblock The Llama 3 herd of models.
\newblock Informe técnico, 2024.
\newblock \url{https://arxiv.org/abs/2407.21783}

\bibitem{mamba}
A.~Gu y T.~Dao.
\newblock Mamba: Linear-time sequence modeling with selective state spaces.
\newblock Prepublicación, 2023 (revisada en 2024).
\newblock \url{https://arxiv.org/abs/2312.00752}

\bibitem{s4}
A.~Gu, K.~Goel y C.~Ré.
\newblock Efficiently modeling long sequences with structured state spaces.
\newblock En \emph{International Conference on Learning Representations (ICLR)}, 2022.
\newblock \url{https://arxiv.org/abs/2111.00396}

\bibitem{cs336}
T.~Hashimoto y P.~Liang.
\newblock CS336: Language Modeling from Scratch.
\newblock Stanford University, primavera de 2025. Clases~3 («Architectures, hyperparameters»), 4 («Mixture of experts») y 10 («Inference»), transparencias y código publicados.
\newblock \url{https://github.com/stanford-cs336/spring2025-lectures}

\bibitem{cs336a1}
Equipo docente de CS336.
\newblock CS336: Language Modeling from Scratch. Tarea~1, «Basics», versión 26.0.3.
\newblock Stanford University, primavera de 2026. §3.4.1 (RMSNorm), §3.4.2 (MLP SwiGLU) y §7.3 (ablaciones de arquitectura).
\newblock \url{https://github.com/stanford-cs336/assignment1-basics}

\bibitem{mixtral}
A.~Q. Jiang \emph{et al.}
\newblock Mixtral of experts.
\newblock Prepublicación, 2024.
\newblock \url{https://arxiv.org/abs/2401.04088}

\bibitem{mistral}
A.~Q. Jiang \emph{et al.}
\newblock Mistral 7B.
\newblock Prepublicación, 2023.
\newblock \url{https://arxiv.org/abs/2310.06825}

\bibitem{jamba}
O.~Lieber \emph{et al.}
\newblock Jamba: A hybrid Transformer-Mamba language model.
\newblock Prepublicación, 2024. Citado a través de la clase~10 de CS336~\cite{cs336}.
\newblock \url{https://arxiv.org/abs/2403.19887}

\bibitem{mit6003}
Massachusetts Institute of Technology.
\newblock 6.003: Signals and Systems.
\newblock MIT OpenCourseWare, otoño de 2011. \emph{Homework}~2, problemas~1 («Finding outputs») y~6 («Unit-sample response»).
\newblock \url{https://ocw.mit.edu/courses/6-003-signals-and-systems-fall-2011/}

\bibitem{yarn}
B.~Peng, J.~Quesnelle, H.~Fan y E.~Shippole.
\newblock YaRN: Efficient context window extension of large language models.
\newblock En \emph{International Conference on Learning Representations (ICLR)}, 2024.
\newblock \url{https://arxiv.org/abs/2309.00071}

\bibitem{whisper}
A.~Radford, J.~W. Kim, T.~Xu, G.~Brockman, C.~McLeavey e I.~Sutskever.
\newblock Robust speech recognition via large-scale weak supervision.
\newblock En \emph{International Conference on Machine Learning (ICML)}, 2023.
\newblock \url{https://arxiv.org/abs/2212.04356}

\bibitem{mqa}
N.~Shazeer.
\newblock Fast Transformer decoding: One write-head is all you need.
\newblock Prepublicación, 2019.
\newblock \url{https://arxiv.org/abs/1911.02150}

\bibitem{shazeerglu}
N.~Shazeer.
\newblock GLU variants improve Transformer.
\newblock Prepublicación, 2020.
\newblock \url{https://arxiv.org/abs/2002.05202}

\bibitem{shazeer2017}
N.~Shazeer, A.~Mirhoseini, K.~Maziarz, A.~Davis, Q.~Le, G.~Hinton y J.~Dean.
\newblock Outrageously large neural networks: The sparsely-gated mixture-of-experts layer.
\newblock En \emph{International Conference on Learning Representations (ICLR)}, 2017.
\newblock \url{https://arxiv.org/abs/1701.06538}

\bibitem{roformer}
J.~Su, Y.~Lu, S.~Pan, A.~Murtadha, B.~Wen e Y.~Liu.
\newblock RoFormer: Enhanced Transformer with rotary position embedding.
\newblock Prepublicación, 2021. Publicado en \emph{Neurocomputing}, 568, 2024.
\newblock \url{https://arxiv.org/abs/2104.09864}

\bibitem{llama2}
H.~Touvron \emph{et al.}
\newblock Llama 2: Open foundation and fine-tuned chat models.
\newblock Prepublicación, 2023.
\newblock \url{https://arxiv.org/abs/2307.09288}

\bibitem{vaswani}
A.~Vaswani, N.~Shazeer, N.~Parmar, J.~Uszkoreit, L.~Jones, A.~N. Gomez, Ł.~Kaiser e I.~Polosukhin.
\newblock Attention is all you need.
\newblock En \emph{Advances in Neural Information Processing Systems 30 (NeurIPS)}, 2017.
\newblock \url{https://arxiv.org/abs/1706.03762}

\bibitem{grok1}
xAI.
\newblock Grok-1: código y pesos publicados (archivos \texttt{README.md}, \texttt{run.py} y \texttt{model.py}).
\newblock 2024. Licencia Apache~2.0.
\newblock \url{https://github.com/xai-org/grok-1}

\bibitem{soundstream}
N.~Zeghidour, A.~Luebs, A.~Omran, J.~Skoglund y M.~Tagliasacchi.
\newblock SoundStream: An end-to-end neural audio codec.
\newblock Prepublicación, 2021.
\newblock \url{https://arxiv.org/abs/2107.03312}

\bibitem{d2l}
A.~Zhang, Z.~C. Lipton, M.~Li y A.~J. Smola.
\newblock \emph{Dive into Deep Learning}, versión 1.0.3.
\newblock Cambridge University Press, 2023.
\newblock Cap.~11, «Attention Mechanisms and Transformers» (§11.6 «Self-Attention and Positional Encoding» y §11.8 «Transformers for Vision»).
\newblock \url{https://d2l.ai/}

\bibitem{rmsnorm}
B.~Zhang y R.~Sennrich.
\newblock Root mean square layer normalization.
\newblock En \emph{Advances in Neural Information Processing Systems 32 (NeurIPS)}, 2019.
\newblock \url{https://arxiv.org/abs/1910.07467}

\end{thebibliography}
